\Needspace{12\baselineskip}
\section{A uniform two-window relation estimate}\label{sec:two-window-anatomy}
The transfer in \cref{thm:transfer} separates four costs: weighted
one-window defects, supports removed at the prime cutoff, shared-prime
graph edges, and the weighted relations estimated here. Only the last
requires the two-window arithmetic. We sum $2^{\rho(x,y)}-1$, not merely
count pairs with a relation, by separating exact rational identities from
residual prime relations in one macroscopic geometry. Word overlaps are
treated separately in \cref{thm:word-field}.

Fix $0<\delta<1$ and $0<\beta_-<\beta_+<\infty$. Throughout this section,
\begin{equation}\label{eq:uniform-band}
 U=U_{M,\delta},\qquad \beta_-\log M\le B=L+1\le\beta_+\log M,
 \qquad Q_B=2^B.
\end{equation}
A factor $M^{o(1)}$ is uniform in this band, with dependence only on the
fixed parameters (and on an explicitly fixed component size when one is
used). All these parameters are fixed before $M$. Every occurrence lies in
\begin{equation}\label{macro:eq:occ-range}
 M^\delta-1\le n\le M+B\le2M.
\end{equation}

\begin{theorem}[Uniform weighted relation profile]\label{thm:master-profile}
For every $\varepsilon>0$,
\begin{align}
 R_{2,\delta}(M,L)
 &:=\sum_{\substack{x,y\in U\\|x-y|>L}}(2^{\rho_L(x,y)}-1)\notag\\
 &\ll_{\delta,\beta_-,\beta_+,\varepsilon}
 M^\varepsilon\left(M^{3/2}Q_B^{1/6}+MQ_B^{2/3}+M^{2/3}Q_B\right)
 \label{eq:master-profile}\\
 &\ll_{\delta,\beta_-,\beta_+,\varepsilon}
 M^\varepsilon\left(M^{5/3}+M^{2/3}Q_B\right).
 \label{eq:master-envelope}
\end{align}
The defining sum is over ordered pairs. In particular
$R_2(N,L)\ll_{C,\varepsilon}N^{5/3+\varepsilon}$ when
$|L-\log_2N|\le C$ in a dyadic block. The same conclusion, with a constant
depending on the fixed upper ratio, holds on $[N,\kappa N)$.
\end{theorem}
The last assertions follow by restriction from \eqref{eq:master-envelope}
at $M=2N$ or $M=\kappa N$ and, for example, $\delta=1/2$; the lower endpoint
is then at least $M^\delta$ for large $N$. No sum of independent block laws
is being taken. The proof of the estimates occupies the rest of this section.

There are three obstacles to counting the relations. Exact rational
identities may create many of them for one identifiable reason; we remove
that subspace first. Residual relations are then controlled by combining
rarity of their hosts with a bound on their dimension. In the terminal
case neither resource is sufficient separately: determinants force
several small kernels into one window, and it is this simultaneous
concentration, rather than the size of the candidate set, that must be
counted. The partition below distinguishes which of these resources is
available.


\subsection{Rational identities and the residual graph}
For a separated pair, use the even-subset description in
\cref{lem:even-subsets}. An exact rational relation can be present before
any primes are exposed. For example, with $x=100$, $y=67$ and $L=3$, the
windows contain the two units $(99,66)$ and $(102,68)$ of ratio $3:2$.
Their product is $6732^2$. Each selected subset in a window has two vertices,
so this gives a nonzero homogeneous relation; its two root occurrences also
make the affine character even. In fact the individual ranks are three and
the stacked rank is five, giving joint probability $1/32$, rather than
$1/64$ for independent starts.

\begin{figure}[htbp]\centering
\begin{tikzpicture}[x=1.5cm,y=1.1cm]
\foreach \i/\t in {0/99,1/100,2/101,3/102}{
 \node[occurrence] (a\i) at (\i,1) {};
 \node[above=2mm of a\i,font=\small] {$\t$};}
\foreach \i/\t in {0/66,1/67,2/68,3/69}{
 \node[occurrence] (b\i) at (\i,0) {};
 \node[below=2mm of b\i,font=\small] {$\t$};}
\draw[eventedge] (a0)--(a1)--(a2);\draw[eventedge] (a1) to[bend left=20] (a3);
\draw[eventedge] (b0)--(b1)--(b2);\draw[eventedge] (b1) to[bend right=20] (b3);
\draw[channel] (a0)--(b0);\draw[channel] (a3)--(b2);
\node[align=left,font=\small,anchor=west] at (3.6,.5)
 {$2\cdot99=3\cdot66$\\$2\cdot102=3\cdot68$};
\end{tikzpicture}
\caption{An exact two-unit relation. The thick edges are rational units;
the thin edges encode the affine start trees. Even subsets of the units
produce relations. This macroscopic example also illustrates why counting
all prime coincidences without quotienting the rational subspace loses
information.}\label{fig:rational-example}
\end{figure}
For coprime $a,b\geq1$, set $h=by-ax$ and
\[
 \mathcal M_{a,b}(h)=\{(i,j)\in I_L^2:ai-bj=h\},
 \qquad m(a,b,h)=|\mathcal M_{a,b}(h)|.
\]
A member of $\mathcal M_{a,b}(h)$ is an exact unit, because
$a(x+i)=b(y+j)$.  Even subsets of the exact units form a subspace
$\mathcal S_{a,b,h}\subset\Rspace(x,y)$.

\begin{lemma}[Rational code and height]\label{macro:lem:rational-code}
If $m=m(a,b,h)\geq1$, then
$\dim\mathcal S_{a,b,h}=m-1$.  If $m\geq2$, then
\[
 \max(a,b)\leq L<B.
\]
This conclusion is independent of the ratio between the endpoints
$x$ and $y$.
\end{lemma}

\begin{proof}
For an exact unit, coprimality gives integers $t$ with
$x+i=bt$ and $y+j=at$.  An even collection $E$ of units satisfies
\[
 \prod_{(i,j)\in E}(x+i)(y+j)
 =(ab)^{|E|}\left(\prod_{(i,j)\in E}t\right)^2,
\]
which is a square.  The even subsets of an $m$-element set have
dimension $m-1$.  Distinct solutions of $ai-bj=h$ differ by
$(b,a)$; since both offsets remain in an interval of diameter $L$,
one has $a,b\leq L$.
\end{proof}

Fix once and for all $A=3$.  Call $(x,y)$ aligned if there are
coprime $a,b\leq B^A$ such that
\begin{equation}\label{macro:eq:aligned}
 |by-ax|<(a+b)B.
\end{equation}

\begin{lemma}[Uniqueness of the macroscopic channel]\label{macro:lem:channel-unique}
For $M$ sufficiently large in terms of
$(\delta,\beta_-,\beta_+)$, an aligned pair has at most one reduced
fraction $(a,b)$ satisfying \eqref{macro:eq:aligned}.  Every channel with at
least two exact units is aligned and therefore has this unique
fraction.
\end{lemma}

\begin{proof}
If $(a,b)$ and $(a',b')$ both satisfy \eqref{macro:eq:aligned}, then
\[
 |(a'b-ab')x|<4B^{2A+1}.
\]
But $x\geq M^\delta$ and $B=O(\log M)$, so the right-hand side is
$o(M^\delta)$.  The integer $a'b-ab'$ is therefore zero for large
$M$, and reducedness gives $(a,b)=(a',b')$.  If a channel has two
units, Lemma~\ref{macro:lem:rational-code} gives $a,b\leq L$, and for any
unit $h=ai-bj$ satisfies $|h|<(a+b)B$.
\end{proof}

When the canonical aligned channel has $m\geq2$, define
\[
 \Srat=\mathcal S_{a,b,h},\qquad \sigma=m-1.
\]
Otherwise define $\Srat=0$ and $\sigma=0$.  In all cases put
\[
 \tau=\rho-\sigma.
\]
Thus no contraction is made when $m_{\rm ex}\leq1$.

\begin{remark}[Disparate scales]\label{macro:rem:disparate}
Suppose, for example, that $x\asymp M^\delta$ and $y\asymp M$.
If there were two exact units, then $a,b\leq L$ and
$x+i=bt$, $y+j=at$, so
\[
 y\leq L(x+B)+B=o(M),
\]
a contradiction.  Hence every genuinely disparate-scale pair has
$m_{\rm ex}\leq1$.  It remains in the ordinary large-prime graph and
is treated by the small-product, shallow-core or deep-core bounds below;
no illegitimate valuation identity or code contraction is used.
\end{remark}

Now form the graph $G_{>B}(x,y)$ on the occurrence set
$\cO=V_x\sqcup V_y$.  For a prime $p>B$, let $\cO_p$ be the
occurrences at which $v_p$ is odd.  Each block has diameter $B-1<p$,
so $|\cO_p|\leq2$, with at most one occurrence in each block.  If
$|\cO_p|=1$, pin that vertex to zero; if $|\cO_p|=2$, join the two
vertices; parallel equations are deduplicated.  A vertex with no
large-prime label is defective.  Let $D$ be the number of defective
vertices and $c$ the number of nontrivial unpinned components.

\begin{lemma}[Resolution and square class of components]\label{macro:lem:graph-resolution}
The solution space of the $p>B$ equations is canonically isomorphic,
after one representative is chosen in each component, to
\[
 \F^D\oplus\F^c.
\]
In particular, $D+2c\leq2B$.  Every nontrivial unpinned component
$C$ meets both blocks and has a unique representation
\[
 \prod_{n\in C}n=d_Cz_C^2,
\]
where $d_C$ is positive, squarefree, and $B$-smooth.
\end{lemma}

\begin{proof}
A pinned component is zero.  All variables in an unpinned component
containing an edge are equal and contribute one free coordinate.  An
isolated defective vertex contributes one free coordinate.  Every
nontrivial component uses at least two occurrences, giving
$D+2c\leq2B$.

Every edge is bipartite.  For each $p>B$, an unpinned component
contains either both odd occurrences of $p$ or none; a single odd
occurrence would pin it.  Hence the total $p$-valuation of the product
of the component is even for every $p>B$.  Its squarefree part is
therefore supported on $p\leq B$, and uniqueness is immediate.
\end{proof}

Assume $m\geq2$.  An exact unit has the form
$x+i=bt$, $y+j=at$ with $a,b<B$.
For every $p>B$, $p\nmid ab$ and therefore
\[
 v_p(x+i)\equiv v_p(t)\equiv v_p(y+j)\pmod2.
\]
If $K_B(t)=1$, the two occurrences are defective; otherwise they form
an isolated unpinned component of size two.  Distinct exact units use
disjoint occurrences.  Let $d$ be the number of exact units of the
first type and $n=m-d$ the number of the second type, and define
\[
 D^\#=D-d,
 \qquad c^\#=c-n,
\]
removing the $n$ exact components from the residual component set
$\Cres$.  If $m<2$, set $D^\#=D$, $c^\#=c$, and let $\Cres$ be all
nontrivial unpinned components.  Notice that $D^\#$ is a dimensional
correction, not literally a number of remaining defective vertices.

\begin{lemma}[Quotient core]\label{macro:lem:quotient-core}
Uniformly for all separated pairs,
\[
 \tau\leq D^\#+c^\#,
 \qquad D^\#+2c^\#\leq2B,
 \qquad D^\#=O(B/\log B).
\]
\end{lemma}

\begin{proof}
For $m\geq2$, the $m$ exact-unit vectors are independent in the
large-prime solution space and the even code has dimension $m-1$.
The parity functional of the first block descends to the quotient and
is nonzero there; hence the even-first-block subspace has dimension
$D+c-m$.  The full relation quotient embeds in that subspace, and the
small-prime equations can only decrease dimension.  Therefore
\[
 \tau\leq D+c-m=(D-d)+(c-n)=D^\#+c^\#.
\]
For $m<2$, the estimate is the direct inclusion
$\Rspace\subset W_{>B}$.  The inequality
$D^\#+2c^\#\leq2B$ follows from $D+2c\leq2B$.  The defect estimate follows from \cref{prop:defects}, applied to each
window on its own dyadic slice; \eqref{eq:uniform-band} makes the constants
uniform.
\end{proof}

From each component of $\Cres$, choose canonically its smallest
large-prime label and then the lexicographically first edge carrying
that label.  This gives
\[
 (i_s,j_s,p_s),\qquad 1\leq s\leq c^\#,
\]
with distinct offsets in each block, distinct primes $p_s>B$, and
$p_s\mid x+i_s$, $p_s\mid y+j_s$.  When $m\geq2$, these cells are not
exact; when $m=1$, only the at most two components meeting the two
occurrences of the unique exact unit can be exceptions.  Put
\[
 P^\#=\prod_{s=1}^{c^\#}p_s,
\]
with the empty product equal to one.


The dimension decomposition gives the exact identity
\begin{equation}\label{eq:weight-split}
 2^\rho-1=(2^\sigma-1)+2^\sigma(2^\tau-1).
\end{equation}
We first sum the rational term. The residual term will be divided into
exhaustive geometric sectors.

\subsection{Hosts and rational mass}
The elementary prime and divisor estimates used below are
\begin{align}
 \pi(B)&\ll B/\log B,&
 \sum_{p\le B}p^{-1/2}+B\sum_{p>B}p^{-3/2}
 &\ll \sqrt B/\log B,\label{macro:eq:prime-half}\\
 \sum_{p>B}p^{-2}&\ll(B\log B)^{-1},&
 \omega(n)+\log\tau(n)&\ll_K\log M/\log\log M
 \quad(1\le n\le M^K).
 \label{macro:eq:omega-tau}
\end{align}
Here $\tau(n)$ in \eqref{macro:eq:omega-tau} is the divisor function,
not the residual dimension. The last bound is used only for fixed $K$.
Partial summation of the prime number theorem proves the prime estimates;
the divisor estimate is standard \cite{Tenenbaum2015,NicolasRobin1983}.
In particular,
\begin{equation}\label{macro:eq:euler-subpoly}
 \prod_{p\le B}(1+c p^{-1/2})=M^{o(1)},\qquad 2^{\pi(B)}=M^{o(1)}
 \quad(c\hbox{ fixed}),
\end{equation}
while a squarefree $B$-smooth integer has size at most
\begin{equation}\label{macro:eq:Bsmooth-height}
 \prod_{p\le B}p=\exp(O(B))=M^{O_{\beta_+}(1)}.
\end{equation}
Let $H_{2,\delta}(M,L)$ count separated ordered pairs in $U^2$ with
$\rho_L(x,y)>0$. We also need $H^{\rm val}_{2,\delta}(M,B)$, the number of
such pairs for which some nonempty subset of $V_x\sqcup V_y$ has square
product, with no parity requirement in either window.
\begin{proposition}[Host bound]\label{macro:prop:host}
Uniformly in the fixed band,
\[
 H_{2,\delta}(M,L)\le H^{\rm val}_{2,\delta}(M,B)\le M^{3/2+o(1)}.
\]
\end{proposition}

\begin{proof}
For each pair counted by $H^{\rm val}_{2,\delta}$, choose a nonempty
square-product subset and its first selected occurrence, ordering first the $x$-block, then the $y$-block, and then
the offsets increasingly.  Suppose first that this occurrence is
$n=x+i$; the other orientation costs a factor two.  Put $K=K_B(n)$.
For every $p\mid K$, the parity of the selected square product forces
an occurrence of odd $p$-valuation in the other block.  It is unique
there because $p>B$ and the block diameter is $B-1$.  Forgetting
compatibility between primes, there are at most $B^{\omega(K)}$
assignments of the primes of $K$ to offsets of the other block.

A fixed assignment determines one residue class for $y$ modulo every
$p\mid K$, hence one class modulo $K$.  The interval $\UM$ has length
less than $M$, so it contains at most $M/K+1$ representatives.  By
\eqref{macro:eq:occ-range}, $K\leq n\leq2M$, and therefore
\[
 M/K+1\leq3M/K,
\]
including $K=1$.  The orientation and selected offset cost at most
$2B$, and hence
\begin{equation}\label{macro:eq:host-raw}
 H^{\rm val}_{2,\delta}(M,B)
 \ll BM\sum_{n\leq2M}\frac{B^{\omega(K_B(n))}}{K_B(n)}.
\end{equation}
Write uniquely $n=a^2ur$, where $u$ is squarefree and $B$-smooth and
$r=K_B(n)$ is squarefree with all prime factors $>B$.  Summing over
$a$ first gives
\begin{align*}
 \sum_{n\leq2M}\frac{B^{\omega(K_B(n))}}{K_B(n)}
 &\leq \sqrt{2M}
 \sum_u u^{-1/2}\sum_r B^{\omega(r)}r^{-3/2}\\
 &\leq \sqrt{2M}
 \prod_{p\leq B}(1+p^{-1/2})
 \prod_{p>B}(1+Bp^{-3/2}).
\end{align*}
For this host estimate, elementary integer sums suffice:
\[
 \sum_{p\le B}p^{-1/2}\le\sum_{n=1}^B n^{-1/2}\le2\sqrt B,
 \qquad B\sum_{p>B}p^{-3/2}
 \le B\int_B^\infty t^{-3/2}\,dt=2\sqrt B.
\]
Since $\log(1+t)\le t$, the weighted kernel sum is at most
$\sqrt{2M}\exp(4\sqrt B)$. Under $B\le\beta_+\log M$,
$B\exp(4\sqrt B)=M^{o(1)}$, uniformly in the length. Substitution in
\eqref{macro:eq:host-raw} proves the claim. The sharper prime-sum bounds
are not needed at this step.
\end{proof}
\begin{proposition}[Rational mass]\label{macro:prop:systematic}
The contributions with canonical height $q=\max(a,b)$ satisfy
\begin{equation}\label{eq:systematic}
 \sum_{q\ge3}(2^\sigma-1)\le MQ_B^{1/3}M^{o(1)},\qquad
 \sum_{q=2}(2^\sigma-1)\le MQ_B^{1/2}M^{o(1)}.
\end{equation}
The geometry-only sum, without translation parameters, satisfies
\begin{equation}\label{macro:eq:geometry-sum}
 \sum_{\substack{(a,b)=1,\ \max(a,b)\ge2,\ h\in\Z\\m(a,b,h)\ge2}}2^{m(a,b,h)-1}\le Q_B^{1/2}M^{o(1)}.
\end{equation}
\end{proposition}
\begin{proof}
At height $q$ there are $O(q)$ reduced pairs $(a,b)$ and $O(qB)$
possible nonempty values of $h$. A unit $(i_0,j_0)$ gives
$x=bt-i_0$, $y=at-j_0$, with $O(M/q)$ choices of $t$.
Units are spaced by $(b,a)$, so $m-1\le B/q$, and $m\ge2$ implies
$q\le L$. Summing $O(MqB2^{B/q})$ for $q\ge3$ gives the first bound;
$q=2$ gives the second. At $q=1$, separation excludes a unit.
Omitting the $O(M/q)$ translation factor and summing the same geometries
proves \eqref{macro:eq:geometry-sum}.
\end{proof}
\begin{remark}[Why the dyadic interval helps, and what it does not remove]\label{rem:dyadic-edge}
The starts $(x,2x)$ cannot both belong to $[N,2N)$. This removes the
volumetric $2:1$ family, not all ratio-two occurrences in two windows:
$(N-1,2N-2)$ can still occur when the starts are $N$ and $2N-1$.
Our master estimate keeps the volumetric branch and therefore also applies
to the finite prefix. In a dyadic restriction, a more precise channel count
improves this particular term, but that improvement is not needed for
\eqref{eq:master-envelope}.
\end{remark}
\subsection{The partition and the shallow sectors}
For a nonaligned pair let $\widetilde k$ be the rank of the small-prime
and two block-parity equations after solving all large-prime equations.
Then
\begin{equation}\label{macro:lem:exact-rank}
 \tau=D^\#+c^\#-\widetilde k.
\end{equation}
Set $T_B=\lfloor(B-2)/3\rfloor$. After removing $2^\sigma-1$ in
\eqref{eq:weight-split}, apply the following tests successively:
\begin{center}\small
\begin{tabularx}{\textwidth}{@{}cX@{}}\toprule
Sector & Defining test, after the preceding tests have failed\\\midrule
1 & $P^\#\le M$.\\
2 & $\sigma>0$ and the canonical channel has $q\le\sqrt{\log B}$.\\
3 & $c^\#\le B/6$.\\
4 & The pair is aligned.\\
5 & $c^\#\le2B/3$.\\
6 & $D^\#\ge3$.\\
7 & Write $c^\#=B-s$; then $s+\widetilde k>T_B$.\\
8 & $s+\widetilde k\le T_B$.\\\bottomrule
\end{tabularx}
\end{center}
The tests are complementary at each step. In sectors 5--8 the pair is
nonaligned, $\sigma=0$, and $D^\#=D$, $c^\#=c$. In sectors 7--8 one
also has $c^\#>2B/3$ and $D^\#\le2$. Thus this is a disjoint partition
of all residual mass, including pairs with very different heights.

\begin{proposition}[Direct bounds in sectors 1--3]\label{prop:shallow-final}
Up to a common factor $M^{o(1)}$, the residual masses of sectors 1, 2 and 3
are bounded, respectively, by
\[
 M^{3/2}+MQ_B^{1/2},\qquad
 M^{3/2}+MQ_B^{1/2}+MQ_B^{1/3},\qquad
 M^{3/2}Q_B^{1/6}.
\]
\end{proposition}
\begin{proof}
In sector 1 the canonical certificate contains exactly $c^\#$ distinct
primes, all greater than $B$. Hence
\begin{equation}\label{eq:growing-moment} 
 (B+1)^{c^\#}\le P^\#\le M,\qquad
 c^\#\le\frac{\log M}{\log(B+1)},\qquad 2^{c^\#}=M^{o(1)}.
\end{equation}
The last estimate is uniform in the fixed positive logarithmic band.
Together with $D^\#=O(B/\log B)$ and
$\tau\le D^\#+c^\#$, it gives $2^\tau=M^{o(1)}$ pointwise.
When $\sigma=0$, only relational hosts contribute, and their number is
$M^{3/2+o(1)}$ by \cref{macro:prop:host}. When $\sigma>0$, use
\[
 2^\sigma(2^\tau-1)
 \le2(2^\sigma-1)M^{o(1)}
\]
and sum the rational mass in \cref{macro:prop:systematic}.
The two branches give
$M^{o(1)}(M^{3/2}+MQ_B^{1/2})$, without a critical-scale substitution
or any restriction on the ratio between the starts.

In sector 2, write the canonical channel as $h=by-ax$.
If a prime $p>B$ joins a nonexact cell $(i,j)$, it divides
$a(x+i)-b(y+j)=ai-bj-h\ne0$. The existence of an exact unit bounds
this integer by $O(qB)$. Thus every residual certificate prime is at
most a fixed multiple of $qB$.
Consequently $c^\#\le\pi(CqB)=o(B)$ uniformly for
$q\le\sqrt{\log B}$. Also $D^\#=o(B)$, so $\tau=o(B)$.
Summing the systematic weights by \eqref{eq:systematic} bounds this
sector. The displayed $M^{3/2}$ term is harmless and also covers a
partition convention that sends a zero-dimensional rational code here.

In sector 3, $\sigma=0$ or $q>\sqrt{\log B}$, hence $\sigma=o(B)$.
The quotient estimate gives
$2^\sigma(2^\tau-1)\le Q_B^{1/6}M^{o(1)}$.
Only relational hosts contribute. Multiply by their $M^{3/2+o(1)}$ count.
\end{proof}

The untreated pairs now have a large certificate product and more than
$B/6$ residual components. The next step either rules out their alignment
or extracts a bounded component whose square class supplies a rigid
Diophantine equation.

\subsection{Diophantine counting and deep cores}
A bounded component forces a product of a fixed number of shifted values
to lie in one square class. We need to count solutions only in a polynomial
height range. Two of the factors already suffice: after localizing their
square classes, they satisfy a Pell equation. This bounded-height argument
keeps the uniformity we need without a height-free integral-point theorem.
\begin{lemma}[Uniform Pell count]\label{macro:lem:pell}
Fix $K_0>0$.  Let $A,C$ be positive squarefree integers with
$A/C\notin\mathbb Q^{\times2}$, and let $e\neq0$.  If
$A,C,|e|\leq M^{K_0}$, then the number of integer solutions of
\[
 Az^2-Cw^2=e,
 \qquad |z|,|w|\leq M^{K_0},
\]
is
\[
 \exp\!\left(O_{K_0}\!\left(\frac{\log M}{\log\log M}\right)\right)
 =M^{o(1)}.
\]
\end{lemma}

\begin{proof}The detailed reduction and uniformity checks are given in \cref{supp:lem:pell}.\end{proof}
\begin{lemma}[Split products in a polynomial range]\label{lem:split-bounded}
Fix $d\ge2$ and $K>0$. Let $h_1,\ldots,h_d$ be distinct integers and
$e$ a positive squarefree integer with $|h_i|,e\le M^K$.
The number of integer pairs $(X,Y)$ satisfying
\[
 |X|\le M^K,\qquad X+h_i>0\ (1\le i\le d),\qquad
              \prod_{i=1}^d(X+h_i)=eY^2
\]
is at most $\exp(O_{d,K}(\log M/\log\log M))$.
\end{lemma}
\begin{proof}
Put $\Delta=\prod_{i<j}(h_i-h_j)\ne0$ and
$S=\{p:p\mid e\Delta\}$. If $p\notin S$, it divides at most one
factor $X+h_i$ and has even valuation in the product. Hence the
squarefree part of every factor is supported on $S$.
In particular
\[
 X+h_1=su^2,\qquad X+h_2=tv^2,
 \qquad su^2-tv^2=h_1-h_2\ne0,
\]
where $s,t$ are positive squarefree integers supported on $S$.
There are at most $4^{|S|}$ choices of $(s,t)$. We need retain only
$s,t\le2M^K$, and $u,v\le(2M^K)^{1/2}$.
For $s\ne t$, their quotient is not a rational square, so
\cref{macro:lem:pell} bounds the solutions uniformly. For $s=t$,
the equation $s(u-v)(u+v)=h_1-h_2$ is bounded by the divisor estimate.
Each pair $(u,v)$ determines $X$, and then at most two values of $Y$.
Finally $|e\Delta|\le M^{O_{d,K}(1)}$, so
$|S|=O_{d,K}(\log M/\log\log M)$; the divisor and Pell bounds have the
same exponential scale. Multiplying the bounds proves the claim.
\end{proof}
This argument counts bounded-height solutions, not all integral points.
The distinction is precisely what makes a single Pell equation sufficient.
The square-class and height bookkeeping is recorded in
\cref{supp:sec:split-localization}.

\begin{lemma}[Normalization of a component]\label{lem:component-normalization}
Let $I,J$ be two nonempty sets of offsets, with $|I|+|J|\le K_0$, and let $d$ be a positive squarefree $B$-smooth integer. If
\begin{equation}\label{eq:component-square}
 P_I(x)Q_J(y)=dz^2,
 \qquad
 P_I(X)=\prod_{i\in I}(X+i),
 \quad
 Q_J(Y)=\prod_{j\in J}(Y+j),
\end{equation}
then, after fixing $x,I,J,d$, there is a positive squarefree integer $e\le M^{O_{K_0}(1)}$ such that
\begin{equation}\label{eq:normalized-mobile}
 Q_J(y)=ev^2.
\end{equation}
The integer $e$ is the squarefree part of $dP_I(x)$.
\end{lemma}

\begin{proof}
In the group $\mathbb Q^\times/\mathbb Q^{\times2}$, the equality \eqref{eq:component-square} gives
\[
 [Q_J(y)]=[dP_I(x)].
\]
The positive squarefree representative of the right-hand class is $e$; the squarefree part of $Q_J(y)$ is therefore $e$. Since $|I|\le K_0$, we have $dP_I(x)\le M^{O_{K_0}(1)}$.
\end{proof}
\begin{lemma}[One-sided count of a component]\label{lem:one-sided-component}
Fix $K_0$. For $x,I,J,d$ fixed as in \cref{lem:component-normalization}, the number of $y\in U_{M,\delta}$ satisfying \eqref{eq:component-square} is
\[
 \begin{cases}
 O(M^{1/2}),&|J|=1,\\
 M^{o_{K_0}(1)},&|J|\ge2.
 \end{cases}
\]
The same assertion holds after exchanging the two blocks.
\end{lemma}

\begin{proof}
\Cref{lem:component-normalization} gives $Q_J(y)=ev^2$.
For $|J|=1$, the equation $y+j=ev^2$ leaves at most
$1+\sqrt{3M/e}=O(M^{1/2})$ values.
For $|J|\ge2$, use \cref{lem:split-bounded}. Its degree is at most
$K_0$, all shifted factors are positive, $y\le M$, and
$e\le M^{O_{K_0}(1)}$. Thus a single fixed polynomial-height exponent
covers every coefficient and variable. The symmetric statement follows
by interchanging the blocks.
\end{proof}
\begin{lemma}[Two singleton components]\label{lem:singleton-pair}
Fix two offsets $i,j$. The number of triples $(x,y,d)$, with $x,y\in U_{M,\delta}$, $d$ positive, squarefree and $B$-smooth, satisfying
\[
 (x+i)(y+j)=dz^2
\]
is
\begin{equation}\label{eq:size2-hosts}
 M\exp\!\left(O\!\left(\frac{\sqrt B}{\log B}\right)\right).
\end{equation}
\end{lemma}

\begin{proof}
Set $X=x+i$, $Y=y+j$. A prime-by-prime analysis of $XY=dz^2$ gives a unique parametrization
\[
 X=ecu^2,
 \qquad
 Y=(d/e)cv^2,
\]
where $e\mid d$, where $c$ is squarefree and $(c,d)=1$. For fixed $e$, the conditions $1\le X,Y\le3M$ force
\[
 c\le C_e:=\frac{3M}{\max(e,d/e)}.
\]
Dropping the ``squarefree'' and ``coprime to $d$'' conditions, which only increases the count, the number of triples $(u,v,c)$ is at most
\begin{align*}
 \sum_{c\le C_e}
 \left(1+\sqrt{\frac{3M}{ec}}\right)
 \left(1+\sqrt{\frac{3Me}{dc}}\right)
 &\ll
 C_e
 +\sqrt{\frac Me}\sqrt{C_e}
 +\sqrt{\frac{Me}d}\sqrt{C_e}
 +\frac M{\sqrt d}\sum_{c\le C_e}\frac1c\\
 &\ll
 \frac{M\log M}{\sqrt d}.
\end{align*}
Here we used $\max(e,d/e)\ge\sqrt d$. After summing over $e\mid d$, the number of solutions for fixed $d$ is
\[
 \ll\frac{M\tau(d)\log M}{\sqrt d}.
\]
Finally
\[
 \sum_{\substack{d\text{ squarefree}\\p\mid d\Rightarrow p\le B}}
 \frac{\tau(d)}{\sqrt d}
 =\prod_{p\le B}\left(1+\frac2{\sqrt p}\right)
 =\exp\!\left(O\!\left(\frac{\sqrt B}{\log B}\right)\right).
\]
The factor $\log M$ is absorbed into the exponential $M^{o(1)}$.
\end{proof}
\begin{lemma}[Hosts of a bounded component]\label{lem:bounded-component}
Fix $K_0$. The number of separated pairs $(x,y)$ whose non-aligned residual graph contains a nontrivial unpinned component meeting both blocks and of size at most $K_0$ is
\begin{equation}\label{eq:bounded-hosts}
 M^{1+o_{K_0}(1)}.
\end{equation}
If this component has exactly two vertices, one in each block, we have the sharper bound
\[
 M\exp\!\left(O\!\left(\frac{\sqrt B}{\log B}\right)\right).
\]
\end{lemma}

\begin{proof}
Let $I,J$ be the offset sets of the component.
\Cref{macro:lem:graph-resolution} provides, with a unique coefficient
$d$, equation \eqref{eq:component-square}. There are only
$B^{O_{K_0}(1)}=M^{o(1)}$ choices of the two offset sets. We separate
three branches.

\paragraph{Two degrees at least two.}
If $|I|,|J|\ge2$, fix $x$, the offsets and $d$.
\Cref{lem:one-sided-component} leaves $M^{o_{K_0}(1)}$ values of
$y$. Summing over the $M$ values of $x$ and the
$2^{\pi(B)}=M^{o(1)}$ coefficients $d$ gives
$M^{1+o_{K_0}(1)}$ hosts.

\paragraph{A single degree equal to one.}
If $|I|=1<|J|$, the same argument fixes $x$ and uses the mobile side of
degree $|J|\ge2$. If $|J|=1<|I|$, we exchange the two blocks. We again
obtain $M^{1+o_{K_0}(1)}$ after summing over $d$.

\paragraph{Two singletons.}
If $|I|=|J|=1$, \cref{lem:singleton-pair} applies directly. It already
counts the triples $(x,y,d)$ and, in particular, performs the
simultaneous summation
\[
 \sum_{d\ B\text{-smooth, squarefree}}\frac{\tau(d)}{\sqrt d}.
\]
No factor $2^{\pi(B)}$ must be reintroduced in this branch.
After the union over the $B^2$ offset pairs, absorbed into the
exponential, the sharp bound remains
\[
 M\exp\!\left(O\!\left(\frac{\sqrt B}{\log B}\right)\right).
\]
The three branches give both assertions.
\end{proof}
\begin{proposition}[No aligned residual core of linear density]\label{macro:prop:aligned}
For every fixed $\alpha>0$, no aligned pair in $\UM^2$ satisfies
$c^\#\geq\alpha B$ once $M$ is sufficiently large in terms of the
fixed parameters and $\alpha$.
\end{proposition}

\begin{proof}
Let $m_{\rm ex}$ be the number of exact units of the canonical
channel.  If $m_{\rm ex}\geq2$, the nondefective exact components have
already been removed in $\Cres$.  If $m_{\rm ex}=1$, make no height or
valuation assumption and remove directly the at most two graph
components containing its two occurrences.  If $m_{\rm ex}=0$, remove
nothing.  In every case, from $c^\#\geq\alpha B$ there remain at least
$\alpha B/2$ nontrivial components containing no exact unit.  Their
total size is at most $2B$, so at least $m\geq\alpha B/4$ have size at
most $K_\alpha=\lceil8/\alpha\rceil$.

For each such component $C_t$ form the column
\[
 \Phi(C_t)=\left(
 (v_p(\prod_{n\in C_t}n)\bmod2)_{p\leq B},
 |C_t\cap V_x|\bmod2,
 |C_t\cap V_y|\bmod2
 \right)
 \in\F^{\pi(B)+2}.
\]
The kernel of the matrix of these $m$ columns has codimension at most
$r=\pi(B)+2$.  Choose
\[
 t=\left\lceil C_\alpha\frac{B}{\log B\log\log B}\right\rceil
\]
with $C_\alpha$ large.  If no nonzero kernel word had weight at most
$2t$, the Hamming bound would give
$\sum_{j\leq t}\binom{m}{j}\leq2^r$.  Since $m\gg_\alpha B$ and
$t=o(m)$,
\[
 \log \binom{m}{t}
 \geq t\log(m/t)-O(t)
 =(C_\alpha+o_\alpha(1))\frac{B}{\log B},
\]
which exceeds $r\log2$ when $C_\alpha$ is fixed sufficiently large.
Thus a nonzero word uses
\[
 O_\alpha\!\left(\frac{B}{\log B\,\log\log B}\right)
\]
components.

Let $I,J$ be the unions of their offsets.  The last two coordinates of
$\Phi$ make $|I|$ and $|J|$ even; the remaining coordinates, together
with the large-prime equations internal to each component, show that
\[
 \prod_{i\in I}(x+i)\prod_{j\in J}(y+j)
\]
is a square.  Put $d_0=|I|+|J|$; then
\[
 d_0\ll_\alpha\frac{B}{\log B\log\log B}.
\]
For the canonical channel $(a,b,h)$ define
\[
 G(T)=\prod_{i\in I}(T+ai)
      \prod_{j\in J}(T+h+bj).
\]
There are no root collisions: one would be an exact unit, and all
components containing such a unit were removed.  At $T=ax$,
\[
 G(ax)=a^{|I|}b^{|J|}
       \prod_{i\in I}(x+i)\prod_{j\in J}(y+j)
\]
is an integer square.  All roots have modulus $O(B^{A+1})$, while
$ax\geq x\geq M^\delta$; in particular $ax\geq2R$ for the root
radius $R=O(B^{A+1})$ once $M$ is large.  \Cref{lem:runge} gives
\[
 ax\leq(Cd_0B^{A+1})^{2d_0}.
\]
The logarithm of the right-hand side is
$O_\alpha(B/\log\log B)=o(B)$, whereas
\[
 \log(ax)\geq\delta\log M
 \geq(\delta/\beta_+)B.
\]
This contradiction closes the aligned deep sector directly at scale
$M$.
\end{proof}
We now work in the nonaligned complement. A component of size at most
$2/\alpha$ exists whenever $c^\#\ge\alpha B$, since the total number of
occurrences is $2B$. The preceding component count therefore gives
$M^{1+o(1)}$ hosts in every fixed positive-density deep sector.
For sector 5, $\tau\le D^\#+c^\#\le2B/3+o(B)$, so its mass is
$MQ_B^{2/3}M^{o(1)}$.
\begin{lemma}[Two defects in one block]\label{macro:lem:two-defects}
The number of starts $x\in\UM$ whose window contains two distinct
$B$-defective occurrences is $M^{o(1)}$.
\end{lemma}

\begin{proof}
Write
$x+i_1=d_1u^2$ and $x+i_2=d_2v^2$, with $d_1,d_2$ squarefree and
$B$-smooth.  If $d_1\neq d_2$, then $d_1/d_2$ is not a square and
\[
 d_1u^2-d_2v^2=i_1-i_2\neq0
\]
is covered by Lemma~\ref{macro:lem:pell}; all coefficients and variables are
in fixed polynomial ranges by \eqref{macro:eq:Bsmooth-height} and
\eqref{macro:eq:occ-range}.  If $d_1=d_2=d$, then $d\mid i_1-i_2$ and
\[
 (u-v)(u+v)=(i_1-i_2)/d,
\]
which has $M^{o(1)}$ solutions by the divisor bound.  There are
$4^{\pi(B)}B^2=M^{o(1)}$ choices of coefficients and offsets.
\end{proof}
\begin{proposition}[Sector 6]\label{macro:prop:sector6}
For nonaligned deep-core pairs with $c^\#>2B/3$ and $D^\#\geq3$,
\[
 R_{\rm res}\leq M^{1/2}\QB M^{o(1)}.
\]
\end{proposition}

\begin{proof}
One block contains at least two defects, so its start has
$M^{o(1)}$ choices by Lemma~\ref{macro:lem:two-defects}.  Since
$c^\#>2B/3$, a residual component has bounded size.  Fix its offsets
and square class at a cost $M^{o(1)}$.  The start in the other block is
then counted by Lemma~\ref{lem:one-sided-component}: at most $M^{1/2}$ choices
when its component degree is one, and $M^{o(1)}$ otherwise.  Hence the
number of hosts is $M^{1/2+o(1)}$.

Finally $D^\#+2c^\#\leq2B$ and
$D^\#=O(B/\log B)$ imply
$D^\#+c^\#\leq B+o(B)$.  Therefore
$2^\tau\leq\QB M^{o(1)}$, and multiplication proves the result.
\end{proof}
For sector 7, integrality gives $s+\widetilde k\ge T_B+1$.
Thus \eqref{macro:lem:exact-rank} and $D^\#\le2$ imply
$\tau\le B+1-T_B=2B/3+O(1)$. The same deep-host count bounds its mass
by $MQ_B^{2/3}M^{o(1)}$.

\subsection{Terminal kernels: count clustering, not just candidates}
It remains to consider sector 8. Write $j=s+\widetilde k\le T_B$.
If a component has size $2+e_t$, then
$\sum_t e_t\le2B-2c^\#=2s$. At most $2s$ components have size greater
than two, so at least $B-3s\ge B-3j\ge2$ are isolated size-two
components. For such a component $\{x+i_t,y+j_t\}$ put
$R_t=K_B(x+i_t)$.
\begin{lemma}[Terminal kernels and determinants]\label{macro:lem:determinant}
For every isolated size-two component,
\[
 R_t=K_B(x+i_t)=K_B(y+j_t)>1.
\]
Kernels belonging to distinct components are coprime.  For $s\neq t$,
\[
 R_sR_t\mid
 \Delta_{s,t}:=(x+i_s)(y+j_t)-(x+i_t)(y+j_s),
\]
and
\[
 0<|\Delta_{s,t}|\leq4MB
\]
for large $M$.
\end{lemma}

\begin{proof}
An isolated unpinned component contains exactly the two occurrences
carrying every one of its $p>B$ labels, so the full kernels agree and
are nontrivial.  A prime cannot label two components, because each
block contains at most one odd occurrence of that prime; hence the
kernels are coprime.  Divisibility of the determinant follows term by
term.

If the determinant vanished, the two occurrence ratios would be equal.
Writing their common reduced ratio as $b/a$ gives
$x+i_s=bu_s$, $y+j_s=au_s$ and similarly with $u_t$.  Subtraction yields
\[
 i_s-i_t=b(u_s-u_t),
 \qquad j_s-j_t=a(u_s-u_t).
\]
Distinct components have distinct offsets, so $u_s\neq u_t$, and
$a,b\leq L$.  The two cells would therefore be two units of an aligned
exact channel, contradicting nonalignment.  Finally
\[
 |\Delta_{s,t}|
 \leq|x+i_s||j_t-j_s|+|y+j_s||i_s-i_t|
 \leq4MB
\]
by \eqref{macro:eq:occ-range}.
\end{proof}
Slice terminal pairs by $X\le\max(x,y)<2X$, where
$M^\delta/2\le X\le M$. For counting, call the larger coordinate $z$
and the smaller one its partner; keep each pair's original canonical data.
The two possible orientations cost only a factor two, and no symmetry
of the canonical selection is assumed. On this slice
$R_sR_t\ll XB$. Set $T_X=(C XB)^{1/2}$ with an absolute constant large
enough for this inequality. At most one isolated kernel exceeds $T_X$.

\begin{lemma}[Partners and a one-dimensional container]\label{lem:terminal-partners}
On a terminal slice, a fixed value of either start has $X^{o(1)}$
possible partners. The number of possible values of either coordinate,
and hence of the larger start, is at most $X^{3/4+o(1)}$.
All constants are uniform in the slice.
\end{lemma}
\begin{proof}
Fix either start and write it as $x$, with partner $y$, for this count
only; the original canonical data of the pair are unchanged.
Choose two isolated components at a cost $O(B^4)=X^{o(1)}$.
Their distinct coprime kernels $R,S>1$ are fixed by $x$. The other
occurrences have the unique forms
$y+j=Ruz^2$, $y+v=Su'w^2$ with independent squarefree $B$-smooth
coefficients $u,u'$.
Thus
\[
 Ruz^2-Su'w^2=j-v\ne0.
\]
The coefficients $Ru,Su'$ are squarefree, of size $O(X)$, and their
quotient is nonsquare because the large-prime supports of $R,S$ are
disjoint. The Pell bound gives $X^{o(1)}$ partners per pair of smooth
coefficients, of which there are $4^{\pi(B)}=X^{o(1)}$.
This argument needs only $y\ll X$, not $y\asymp X$.

At least one isolated kernel is at most $T_X$. Integers with such a
kernel have the unique representation $n=ua^2r$, where $u$ is squarefree
and $B$-smooth and $r=K_B(n)\le T_X$. Therefore
\begin{equation}\label{eq:kernel-container}
 \#\{n\le3X:K_B(n)\le T_X\}
 \le\sqrt{3X}\prod_{p\le B}(1+p^{-1/2})\sum_{r\le T_X}r^{-1/2}
 \ll X^{3/4+o(1)}.
\end{equation}
Each such integer lies in at most $B=X^{o(1)}$ windows, whether the
coordinate is the first or the second start. Their union also bounds the
larger starts. Since
$B/\log X$ stays in a fixed compact interval, all constants are uniform.
\end{proof}

\begin{lemma}[Shifted small-kernel energy]\label{lem:kernel-energy}
For $T\ll(XB)^{1/2}$ put
$A_T(x)=\#\{i\in I_L:K_B(x+i)\le T\}$. Uniformly on every slice,
\begin{equation}\label{eq:kernel-energy}
 \sum_{X\le x<2X}\binom{A_T(x)}2\ll_\varepsilon X^{2/3+\varepsilon}B^2.
\end{equation}
\end{lemma}
\begin{proof}
Let $W_B=\prod_{p\le B}(1+p^{-1/2})=X^{o(1)}$.
Kernels $r\le R_0=X^{1/3}$ occur at at most
\[
 \sqrt{3X}W_B\sum_{r\le R_0}r^{-1/2}=X^{2/3+o(1)}
\]
integers in $[X/2,3X]$. Each can anchor at most $O(B^2)$ pairs of
occurrences and windows, so these integers give the required bound.

For the remaining pairs fix $1\le h\le B$ and dyadic ranges
$r=K_B(n)\asymp R$, $s=K_B(n+h)\asymp S$, with
$R,S\ge X^{1/3}$ and $R,S\ll(XB)^{1/2}$.
Set $q=n/r$, $q'=(n+h)/s$. Both are a square times a squarefree
$B$-smooth integer. There are at most
$X^{o(1)}\sqrt{X/R}$ and $X^{o(1)}\sqrt{X/S}$ choices, respectively.
Moreover $q\asymp X/R$ and $q'\asymp X/S$, because $n\asymp X$.
For a fixed pair, the equation
\[
 qr+h=q's
\]
requires $(q,q')\mid h$; if soluble, its integer solutions have steps
$q'/(q,q')$ and $q/(q,q')$. Hence at most $1+O(BRS/X)$ solutions lie
in the specified boxes. Their total is
\[
 X^{o(1)}\left(\frac{X}{\sqrt{RS}}+B\sqrt{RS}\right)
 \le X^{2/3+o(1)}.
\]
There are $X^{o(1)}$ pairs of dyadic ranges. Sum over $h$ and note
that a shifted pair lies in at most $B$ windows. This proves
\eqref{eq:kernel-energy}.
\end{proof}

\begin{proposition}[Terminal mass]\label{prop:terminal-final}
Sector 8 contributes at most
\[
 M^\varepsilon\left(M^{2/3}Q_B+M^{3/4}Q_B^{2/3}\right).
\]
\end{proposition}
\begin{proof}
On a fixed larger-start slice put $j=s+\widetilde k$ and
$m_j=B-3j-1$. The sector test gives $m_j\ge1$.
By \cref{macro:lem:exact-rank},
$\tau+j=B+D^\#$ with $D^\#\le2$.
There are at least $B-3j$ isolated kernels, all but at most one below
$T_X$, so each window has at least $m_j$ small kernels.

First collect all pairs with $m_j\ge2$, without separating the values
of $j$. Their larger start $z$ satisfies
$\binom{A_{T_X}(z)}2\ge1$. The energy estimate counts at most
$X^{2/3+o(1)}B^2$ such starts, and the uniform partner bound contributes
only $X^{o(1)}$. Since $2^\tau-1\le4Q_B$, the total mass is at most
$X^{2/3+o(1)}Q_B$; the factor $B^2$ is subpolynomial.

In the complementary population, $m_j\le1$ implies $3j\ge B-2$.
Consequently $3\tau\le2B+8$, and
$2^\tau-1\le8Q_B^{2/3}$. At least one small kernel remains, so the
one-dimensional container and partner bound give at most
$X^{3/4+o(1)}$ pairs. Their mass is at most
$X^{3/4+o(1)}Q_B^{2/3}$.

The dyadic slices of the larger coordinate exhaust the terminal pairs.
Their logarithmic bands are uniform, and geometric summation proves the
claim. The partner estimate uses only its upper height bound, never
comparability with $X$.
\end{proof}

\subsection{Assembly and lower bounds}
\begin{proof}[Proof of \cref{thm:master-profile}]
The completed estimates are as follows; a factor $M^{o(1)}$ is suppressed.
\begin{center}\small
\begin{tabularx}{\textwidth}{@{}lXl@{}}\toprule
Population & Mechanism & Weighted mass\\\midrule
Rational & Exact-unit enumeration & $MQ_B^{1/3}+MQ_B^{1/2}$\\
1 & Small product and rational mass & $M^{3/2}+MQ_B^{1/2}$\\
2 & Small channel height & $M^{3/2}+MQ_B^{1/2}+MQ_B^{1/3}$\\
3 & Host count and rank & $M^{3/2}Q_B^{1/6}$\\
4 & Runge and a short codeword & $0$\\
5 & Bounded components & $MQ_B^{2/3}$\\
6 & Two defects and a component & $M^{1/2}Q_B$\\
7 & Exact rank deficit & $MQ_B^{2/3}$\\
8 & Shifted kernel energy and Pell & $M^{2/3}Q_B+M^{3/4}Q_B^{2/3}$\\\bottomrule
\end{tabularx}
\end{center}
Every entry is bounded by one of the three monomials in
\eqref{eq:master-profile}. The partition is exhaustive and
\eqref{eq:weight-split} is exact, so summation proves that profile.
Finally
\[
 M^{3/2}Q_B^{1/6}=(M^{5/3})^{5/6}(M^{2/3}Q_B)^{1/6},\qquad
 MQ_B^{2/3}=(M^{5/3})^{1/3}(M^{2/3}Q_B)^{2/3};
\]
weighted arithmetic--geometric mean proves \eqref{eq:master-envelope}.
\end{proof}

The exponent $5/3$ is an upper bound, not an optimality claim. In a dyadic
block, the ordered pairs $x=2t,y=3t$, $N/2\le t<2N/3$, and their reversals
give
\begin{equation}\label{eq:rational-lower-dyadic}
 R_2(N,L)\ge\left(\frac N3+O(1)\right)
       \left(2^{\lfloor(B+1)/3\rfloor-1}-1\right).
\end{equation}
On $U_{M,\delta}$, the pairs $x=t,y=2t$ and their reversals give
\begin{equation}\label{eq:rational-lower-macro}
 R_{2,\delta}(M,L)\ge(M-2M^\delta+O(1))
       \left(2^{\lfloor B/2\rfloor-1}-1\right).
\end{equation}
These follow simply by counting the even code of exact units, regardless
of any extra residual relations. At criticality the lower exponents are
$4/3$ and $3/2$, respectively. The one-dimensional container
\eqref{eq:kernel-container} is itself of order $X^{3/4+o(1)}$ when
$T\asymp\sqrt{XB}$: choose $r$ prime near $T$ and $n=ra^2\asymp X$.
What improves the terminal estimate is not a smaller container, but the
fact that its elements cannot cluster sufficiently often in one short
window, as quantified by \eqref{eq:kernel-energy}.


The weighted relation estimate is now available. A Poisson comparison
still requires the conditional marginal probabilities, an exact dependency
graph and the local overlap counts. Before constructing that graph, we
record two forms of the estimate adapted to prescribed words.

\subsection{From relative signs to prescribed values}\label{sec:absolute-relations}
There is little extra arithmetic in prescribing the sign of a word rather
than identifying it with its negative. What changes is the parity
restriction on square-product subsets.

For separated windows let $\rho^{\rm val}_B(x,y)$ be the nullity of the
transpose of the matrix with all $2B$ valuation rows $v(n)$,
$n\in V_x\sqcup V_y$. Thus its kernel consists of all square-product
subsets, without requiring even cardinality in each block. Write
$R^{\rm val}_{2,\delta}(M,B)$ for the sum of
$2^{\rho^{\rm val}_B(x,y)}-1$ over the same separated pairs as $R_{2,\delta}$,
and $R^{\rm val}_2(N,B)$ for its dyadic analogue.

\begin{proposition}[Prescribing the global signs]\label{prop:absolute-profile}
Uniformly under the hypotheses of \cref{thm:master-profile},
\begin{align}
 R^{\rm val}_{2,\delta}(M,B)
 &\le4R_{2,\delta}(M,B-1)+3H^{\rm val}_{2,\delta}(M,B),
       \label{eq:absolute-profile-comparison}\\
 R^{\rm val}_{2,\delta}(M,B)
 &\ll_\varepsilon M^\varepsilon
   \left(M^{3/2}Q_B^{1/6}+MQ_B^{2/3}+M^{2/3}Q_B\right)
 \ll_\varepsilon M^\varepsilon(M^{5/3}+M^{2/3}Q_B).
       \label{eq:absolute-profile}
\end{align}
The same conclusion holds on a dyadic block and on bounded-ratio
subintervals, with the corresponding ambient scale.
\end{proposition}
\begin{proof}
Imposing the two block-parity equations on the full square-relation
kernel gives exactly the relative-sign kernel, by the tree boundary
bijection. Therefore
\[
 \rho_{B-1}(x,y)\le\rho^{\rm val}_B(x,y)
                         \le\rho_{B-1}(x,y)+2.
\]
If the full kernel is zero both weights vanish. Otherwise
\[
 2^{\rho^{\rm val}_B}-1
 \le4(2^{\rho_{B-1}}-1)+3.
\]
The host term includes every separated ordered pair admitting a
nonempty square-product subset of the full vertex occurrences, with no
even-cardinality condition in either block. In particular it may count a
pair whose relative-sign kernel is zero. Sum the inequality using exactly
this unrestricted host estimate from \cref{macro:prop:host}. Its $M^{3/2+o(1)}$ term is absorbed by the first
monomial in \eqref{eq:master-profile}. Restriction proves the other
geometries. The extra two affine constraints therefore cost only a fixed
factor and a sparse set of new relational pairs, not an exponential
factor in the window length.
\end{proof}


\subsection{Capping the relation weight}\label{sec:capped-relations}
The full relation weight is useful when estimating moments. For the
probability that one of many mutually exclusive words occurs, however,
its largest values can be wasteful: a joint probability cannot exceed
either marginal. We record a version of the arithmetic estimate that
preserves this elementary ceiling.

For $T\ge1$ define
\[
 R_{2,\delta}^{[T]}(M,L)
   =\sum_{\substack{x,y\in U_{M,\delta}\\|x-y|>L}}
                         \min\{T,2^{\rho_L(x,y)}-1\}.
\]
The notation $R_{2,\delta}^{\mathrm{val},[T]}(M,B)$ has the same meaning
with the full-value nullity $\rho_B^{\rm val}$; omission of $\delta$
denotes the corresponding dyadic sum.
\begin{proposition}[Capped two-window profile]\label{prop:capped-profile}
Under the fixed-band hypotheses of \cref{thm:master-profile}, uniformly
also for $T\ge1$,
\begin{align}
 R_{2,\delta}^{[T]}(M,L)
 &\ll_\varepsilon M^\varepsilon
 \left(M^{3/2}Q_B^{1/6}+MQ_B^{2/3}
                    +M^{2/3}\min\{T,Q_B\}\right),
       \label{eq:capped-profile}\\
 R_{2,\delta}^{\mathrm{val},[T]}(M,B)
 &\ll_\varepsilon M^\varepsilon
 \left(M^{3/2}Q_B^{1/6}+MQ_B^{2/3}
                    +M^{2/3}\min\{T,Q_B\}\right).
       \label{eq:capped-value-profile}
\end{align}
The same bounds hold on dyadic and bounded-ratio subintervals. The
implicit constants are independent of $T$.
\end{proposition}
\begin{proof}
Write $2^\rho-1=(2^\sigma-1)+2^\sigma(2^\tau-1)$ and use
$\min(T,a+b)\le a+\min(T,b)$ for $a,b\ge0$.
The rational contribution and sectors 1--5 and 7 are left uncapped;
the assembly table bounds their sum by
$M^{o(1)}(M^{3/2}Q_B^{1/6}+MQ_B^{2/3})$.
In sector 6, the proof of its host bound gives $M^{1/2+o(1)}$ pairs,
and the pointwise weight is at most $Q_BM^{o(1)}$. Here the factor is
chosen independently of $T$. The inequality
$\min(T,EQ_B)\le E\min(T,Q_B)$ for $E\ge1$ therefore bounds this
contribution by $M^{1/2+o(1)}\min(T,Q_B)$.

For sector 8, use the two populations in the proof of
\cref{prop:terminal-final}. On a larger-start slice, the first population
has at most $X^{2/3+o(1)}$ pairs and pointwise weight at most $4Q_B$.
All geometric and partner bounds were fixed before $T$, so
\[
 \min\{T,2^\tau-1\}\le4\min\{T,Q_B\}
\]
gives capped mass $X^{2/3+o(1)}\min\{T,Q_B\}$ with a common
subpolynomial factor independent of $T$. Leave the complementary
population uncapped: its mass
$X^{3/4+o(1)}Q_B^{2/3}$ is absorbed by $X^{1+o(1)}Q_B^{2/3}$.
Geometric summation proves \eqref{eq:capped-profile}.

Finally, the two parity constraints give the pointwise inequality
\[
 \min\{T,2^{\rho_B^{\rm val}}-1\}
 \le4\min\{T,2^{\rho_{B-1}}-1\}
                       +3\1_{\{\rho_B^{\rm val}>0\}}.
\]
The last indicator has sum $M^{3/2+o(1)}$ by the full-value host
bound. This proves \eqref{eq:capped-value-profile}; restriction proves
the other geometries.
\end{proof}
The cap changes only the dense terminal contribution. Its benefit will
come from choosing $T$ according to a probability, rather than from
improving the uncapped exponent $5/3$.
